\chapter*{Abstract}

\section*{\tituloPortadaEngVal}

Video games are currently a highly widespread and normalized form of entertainment. Unfortunately, since they usually require hand-based controls, there are people who are unable to play them due to mobility limitations. Solutions such as the \textit{QuadStick} exist, allowing these users to play using their mouth. However, these systems rely on controls that are often unintuitive and impractical. An ideal alternative would be enabling users to perform actions in a game simply by thinking about what they want to do. To achieve this, attention must be focused on the brain itself.

The acquisition of brain waves through electroencephalography (EEG) has existed for many years, and its recent integration with artificial intelligence (AI) models has opened new possibilities for their interpretation. To date, brain-computer interface (BCI) systems have achieved reliable results mainly through reactive paradigms that depend on constant external stimuli, which is not suitable for dynamic systems such as video games. Only a limited number of studies have obtained promising results without external stimuli, using motor imagery (MI), and those that do generally rely on hardware that is not easily accessible due to its cost or limited portability.

This work proposes the design and implementation of a modular BCI system capable of processing EEG signals in real time and translating them into commands for controlling video games. The system employs a lightweight and affordable electroencephalograph together with the deep learning model MiRepNet, adapted through \textit{fine-tuning} techniques to capture the specific characteristics of each user.

Experimental results show an average accuracy of 80\% in binary motor imagery classification (right hand versus left hand) on recorded (\textit{offline}) data, reaching values above 90\% for subjects with high-quality signals.

To validate the system in a dynamic context, two video games (\textit{Endless Wave Runner} and \textit{Project Arrows}) were developed, integrating decision filters and assistance mechanisms or \textit{``aimbots''}. These mechanisms compensate for the variability of brain signals, providing a smooth and enjoyable user experience.

Once the complete system had been designed, implemented, and assembled, it was tested experimentally. The experiments demonstrated that mental control of complex applications is possible using accessible hardware. Most of the users who tested the system reported being satisfied with the experience, describing it as both fun and entertaining to play using only the mind.

\section*{Keywords}

\noindent BCI, motor imagery (MI), artificial intelligence, EEG, video games, real-time, signal processing, accessibility, MiRepNet, human-computer interaction.